\section{Fibracje symplicjalne i kategorie modelowe}\label{cat:sec:simplicial-model-categories}

W rozdziale~\ref{cat:sec:simplicial-sets} rozważaliśmy przedłużanie morfizmów z rogów sympleksów, a w rozdziale~\ref{cat:sec:quasicategories} zastosowaliśmy tę własność do quasi-kategorii. Teraz przejdziemy od własności obiektów do własności morfizmów, a następnie do abstrakcyjnej teorii kategorii modelowych. Fibracje topologiczne i ich długie ciągi dokładne należą do topologii algebraicznej; nie powtarzamy tutaj ich konstrukcji.

\subsection{Własności podnoszenia w dowolnej kategorii}

\begin{definicja}[Własność podnoszenia]\label{cat:def:lifting-property}
Niech $i:A\to B$ i $p:E\to Y$ będą morfizmami w kategorii $\mathcal C$. Piszemy $i\pitchfork p$, jeżeli dla każdej pary $a:A\to E$, $b:B\to Y$ spełniającej $pa=bi$ istnieje $l:B\to E$ taki, że
\[
 li=a,\qquad pl=b.
\]
Mówimy, że $i$ ma lewą własność podnoszenia względem $p$, albo równoważnie, że $p$ ma prawą własność podnoszenia względem $i$. Nie wymagamy jednoznaczności $l$.
\end{definicja}

\begin{center}
\begin{tikzpicture}[>=stealth]
 \node (A) at (0,2.6) {$A$};
 \node (E) at (4.4,2.6) {$E$};
 \node (B) at (0,0) {$B$};
 \node (Y) at (4.4,0) {$Y$};
 \draw[->] (A)--node[above] {$a$}(E);
 \draw[->] (A)--node[left] {$i$}(B);
 \draw[->] (E)--node[right] {$p$}(Y);
 \draw[->] (B)--node[below] {$b$}(Y);
 \draw[->,dashed] (B)--node[above,sloped] {$l$}(E);
\end{tikzpicture}
\end{center}

Dla klas morfizmów $\mathcal L,\mathcal R$ oznaczamy
\[
 {}^\pitchfork\mathcal R=\{i:\ \forall p\in\mathcal R\quad i\pitchfork p\},
 \qquad
 \mathcal L^\pitchfork=\{p:\ \forall i\in\mathcal L\quad i\pitchfork p\}.
\]
W szczególności $\mathcal L\subseteq{}^\pitchfork(\mathcal L^\pitchfork)$ oraz $\mathcal R\subseteq({}^\pitchfork\mathcal R)^\pitchfork$. Obie operacje odwracają inkluzje klas, ponieważ zwiększenie klasy morfizmów, względem których wymagamy podnoszenia, zwiększa liczbę warunków do spełnienia.

\begin{stwierdzenie}[Składanie morfizmów z własnością podnoszenia]\label{cat:prop:lifting-composition}
Klasy $\mathcal L^\pitchfork$ i $\vphantom{R}^\pitchfork\mathcal R$ zawierają wszystkie izomorfizmy i są domknięte względem składania.
\end{stwierdzenie}
\begin{proof}
Dla izomorfizmu $i$ w kwadracie z definicji~\ref{cat:def:lifting-property} bierzemy $l=ai^{-1}$: mamy $li=a$ oraz $pl=pai^{-1}=bii^{-1}=b$. Dla izomorfizmu $p$ bierzemy $l=p^{-1}b$, co daje $pl=b$ i $li=p^{-1}bi=p^{-1}pa=a$.

Niech $i:A\to B$, $j:B\to C$ mają lewą własność podnoszenia względem $p:E\to Y$. Dany kwadrat z $ji$ ma dane $a:A\to E$ i $b:C\to Y$ oraz $pa=bji$. Najpierw stosujemy $i\pitchfork p$ do $a,bj$, otrzymując $u:B\to E$ z $ui=a$, $pu=bj$. Następnie $j\pitchfork p$ daje $l:C\to E$ z $lj=u$, $pl=b$. Stąd $lji=ui=a$.

Dla $p:E\to Y$, $q:Y\to Z$ mających prawą własność podnoszenia względem $i:A\to B$ i danych $a:A\to E$, $b:B\to Z$ z $qpa=bi$ najpierw stosujemy $i\pitchfork q$ do $pa,b$. Otrzymujemy $v:B\to Y$ z $vi=pa$, $qv=b$. Stosując $i\pitchfork p$ do $a,v$, uzyskujemy $l:B\to E$ z $li=a$, $pl=v$, a więc $qpl=b$. Dowody działają dla każdego morfizmu z ustalonej klasy.
\end{proof}

\begin{stwierdzenie}[Zmiana bazy i kobazy]\label{cat:prop:lifting-base-change}
Klasa $\mathcal L^\pitchfork$ jest domknięta względem istniejących zmian bazy (iloczynów włóknistych), a klasa $\vphantom{R}^\pitchfork\mathcal R$ względem istniejących zmian kobazy (sum amalgamowanych).
\end{stwierdzenie}
\begin{proof}
Niech $p:E\to Y$ ma prawą własność podnoszenia względem $i:A\to B$, a $f:Y'\to Y$ będzie dowolnym morfizmem. Piszemy $E'=Y'\times_Y E$ z projekcjami $p':E'\to Y'$ i $r:E'\to E$. Dla kwadratu $p'a=bi$ z $a:A\to E'$ i $b:B\to Y'$ mamy $p(ra)=fp'a=fbi$. Z $i\pitchfork p$ dostajemy $t:B\to E$ z $ti=ra$, $pt=fb$. Własność uniwersalna iloczynu włóknistego daje $l:B\to E'$ z $p'l=b$ i $rl=t$. Ponieważ $p'li=bi=p'a$ oraz $rli=ti=ra$, jednoznaczność w tej własności uniwersalnej daje $li=a$.

Dla drugiej tezy niech $i:A\to B$ ma lewą własność podnoszenia względem $p:E\to Y$. Weźmy sumę amalgamowaną $B'=B\amalg_A A'$ względem $f:A\to A'$ z morfizmami $r:B\to B'$ i $i':A'\to B'$. Dla $a:A'\to E$, $b:B'\to Y$ takich, że $pa=bi'$, mamy $paf=bi'f=bri$. Istnieje więc $t:B\to E$ z $ti=af$ i $pt=br$. Z własności uniwersalnej sumy amalgamowanej dostajemy $l:B'\to E$ z $li'=a$ i $lr=t$. Morfizmy $pl,b$ zgadzają się po prekompozycji z $r$ i $i'$, zatem $pl=b$.
\end{proof}

\begin{definicja}[Retrakt morfizmu]\label{cat:def:arrow-retract}
Morfizm $i:A\to B$ jest retraktem morfizmu $j:C\to D$, jeśli istnieją morfizmy $u:A\to C$, $v:B\to D$, $r:C\to A$, $s:D\to B$ takie, że
\[
 ju=vi,\qquad ir=sj,\qquad ru=\id_A,\qquad sv=\id_B.
\]
Jest to zwykły retrakt obiektu w kategorii morfizmów $\operatorname{Fun}([1],\mathcal C)$.
\end{definicja}

\begin{stwierdzenie}[Domknięcie względem retraktów]\label{cat:prop:lifting-retract}
Obie klasy $\mathcal L^\pitchfork$ oraz $\vphantom{R}^\pitchfork\mathcal R$ są domknięte względem retraktów morfizmów.
\end{stwierdzenie}
\begin{proof}
Niech $i$ będzie retraktem $j$ jak w definicji~\ref{cat:def:arrow-retract} i niech $j\pitchfork p$. Dla kwadratu $pa=bi$ mamy $par=bir=bsj$. Istnieje $t:D\to E$ z $tj=ar$, $pt=bs$. Definiujemy $l=tv$. Wtedy
\[
 li=tvi=tju=aru=a,\qquad pl=ptv=bsv=b.
\]
Zatem $i\pitchfork p$.

Dla prawej własności niech $p:E\to Y$ będzie retraktem $q:E'\to Y'$ z morfizmami $u:E\to E'$, $v:Y\to Y'$, $r:E'\to E$, $s:Y'\to Y$ spełniającymi $qu=vp$, $pr=sq$, $ru=\id_E$, $sv=\id_Y$. Jeżeli $i\pitchfork q$ i $pa=bi$, to $qua=vbi$. Istnieje $t:B\to E'$ z $ti=ua$ i $qt=vb$. Dla $l=rt$ mamy $li=rua=a$ oraz $pl=sqt=svb=b$. Zatem $i\pitchfork p$.
\end{proof}

\subsection{Fibracje zbiorów symplicjalnych}

W dalszej części oznaczamy przez $\mathbf{sSet}$ kategorię zbiorów symplicjalnych. Jej obiekty i morfizmy opisano w rozdziale~\ref{cat:sec:simplicial-sets}. Granice i kogranice małych diagramów w tej kategorii obliczamy stopniowo w $\Set$: operatory symplicjalne powstają z własności uniwersalnych odpowiednich granic lub kogranic zbiorów. Tożsamości operatorów wynikają z jednoznaczności w tych własnościach uniwersalnych. W szczególności morfizm $X\to Y$ jest monomorfizmem wtedy i tylko wtedy, gdy wszystkie jego składowe $X_n\to Y_n$ są injekcjami. Wystarczalność wynika ze stopniowej równości morfizmów, a konieczność z testowania monomorfizmu na morfizmach $\Delta[n]\to X$, odpowiadających elementom $X_n$ przez lemat Yonedy (stwierdzenie~\ref{cat:prop:simplicial-yoneda}).

\begin{definicja}[Fibracje symplicjalne]\label{cat:def:simplicial-fibrations}
Morfizm $p:X\to Y$ zbiorów symplicjalnych nazywamy:
\begin{enumerate}
 \item fibracją Kana, jeśli ma prawą własność podnoszenia względem wszystkich włożeń $\Lambda^k[n]\to\Delta[n]$ dla $n\geq1$, $0\leq k\leq n$;
 \item fibracją wewnętrzną, jeśli ma tę własność dla $n\geq2$, $0<k<n$;
 \item lewą fibracją, jeśli ma tę własność dla $n\geq1$, $0\leq k<n$;
 \item prawą fibracją, jeśli ma tę własność dla $n\geq1$, $0<k\leq n$.
\end{enumerate}
\end{definicja}
Każda fibracja Kana jest lewą i prawą fibracją, a każda lewa lub prawa fibracja jest fibracją wewnętrzną. Morfizm jest fibracją Kana wtedy i tylko wtedy, gdy jest zarazem lewą i prawą fibracją, ponieważ sumą obu rodzin indeksów jest rodzina wszystkich rogów.

\begin{center}
\begin{tikzpicture}[>=stealth]
 \node (A) at (0,2.8) {$\Lambda^k[n]$};
 \node (E) at (4.8,2.8) {$X$};
 \node (B) at (0,0) {$\Delta[n]$};
 \node (Y) at (4.8,0) {$Y$};
 \draw[->] (A)--node[above] {$a$}(E);
 \draw[->,hook] (A)--(B);
 \draw[->] (E)--node[right] {$p$}(Y);
 \draw[->] (B)--node[below] {$b$}(Y);
 \draw[->,dashed] (B)--node[above,sloped] {$l$}(E);
\end{tikzpicture}
\end{center}
Różnica względem warunku na obiekt polega na tym, że nie wybieramy dowolnego wypełnienia rogu w $X$, lecz takie, którego obraz przez $p$ jest zadanym sympleksem $b$.

\begin{stwierdzenie}[Postać stopniowa warunku podnoszenia]\label{cat:prop:relative-horn-condition}
Prawa własność podnoszenia $p:X\to Y$ względem $\Lambda^k[n]\to\Delta[n]$ jest równoważna następującemu warunkowi: dla każdego $y\in Y_n$ i każdej rodziny $x_i\in X_{n-1}$, $i\ne k$, spełniającej
\[
 d_i x_j=d_{j-1}x_i\quad(i<j,\ i,j\ne k),\qquad
 p_{n-1}(x_i)=d_i y\quad(i\ne k),
\]
istnieje $x\in X_n$ z $d_i x=x_i$ dla $i\ne k$ i $p_n(x)=y$. Dla $n=1$ pierwszy warunek jest pusty.
\end{stwierdzenie}
\begin{proof}
Morfizm $\Delta[n]\to Y$ odpowiada $y\in Y_n$. Morfizm $\Lambda^k[n]\to X$ odpowiada rodzinie zgodnych ścian na mocy stwierdzenia~\ref{cat:prop:simplicial-compatible-faces}. Przemienność kwadratu podnoszenia oznacza, że dwa morfizmy z rogu do $Y$ są równe. Ponieważ róg jest sumą zadanych ścian, wystarczy i trzeba sprawdzić równość ich ograniczeń, czyli $p_{n-1}(x_i)=d_i y$. Morfizm podnoszący odpowiada $x\in X_n$. Przemienność dwóch trójkątów oznacza dokładnie pozostałe równości.
\end{proof}

\begin{stwierdzenie}[Włókna i obiekty]\label{cat:prop:simplicial-fibration-fibers}
Niech $y\in Y_0$. Włókno $X_y=X\times_Y\{y\}$ fibracji Kana jest kompleksem Kana, a włókno fibracji wewnętrznej jest quasi-kategorią. Morfizm $X\to\Delta[0]$ jest fibracją Kana wtedy i tylko wtedy, gdy $X$ jest kompleksem Kana, oraz fibracją wewnętrzną wtedy i tylko wtedy, gdy $X$ jest quasi-kategorią. Jeżeli $p:X\to Y$ jest fibracją wewnętrzną i $Y$ jest quasi-kategorią, to $X$ jest quasi-kategorią; analogicznie dla fibracji i kompleksów Kana.
\end{stwierdzenie}
\begin{proof}
Morfizm do $\Delta[0]$ jest jedyny, ponieważ w każdym stopniu $\Delta[0]$ ma dokładnie jeden element. Każdy kwadrat podnoszenia dla $X\to\Delta[0]$ jest więc po prostu morfizmem z rogu do $X$, a morfizm podnoszący jest jego wypełnieniem. Daje to obie charakteryzacje.

Projekcja $X_y\to\{y\}$ ma odpowiednią własność podnoszenia przez stwierdzenie~\ref{cat:prop:lifting-base-change}. Ponieważ $\{y\}\cong\Delta[0]$, stosujemy powyższe charakteryzacje. Ostatnie zdania wynikają z domknięcia względem składania (stwierdzenie~\ref{cat:prop:lifting-composition}), zastosowanego do $X\to Y\to\Delta[0]$.
\end{proof}

\begin{przyklad}[Rzuty]\label{cat:ex:simplicial-product-fibration}
Jeżeli $K$ jest kompleksem Kana, to dla dowolnego $Y$ rzut $Y\times K\to Y$ jest fibracją Kana. Jest to zmiana bazy morfizmu $K\to\Delta[0]$ przez $Y\to\Delta[0]$. Gdy $K$ jest quasi-kategorią, ten sam argument daje fibrację wewnętrzną. W szczególności dowolny iloczyn małej rodziny kompleksów Kana jest kompleksem Kana: dla rogu w iloczynie wybieramy wypełnienie każdej składowej i bierzemy odpowiadający im sympleks iloczynu. Korzystamy tu z aksjomatu wyboru.
\end{przyklad}

\subsection{Fibracje na nerwach kategorii}

\begin{stwierdzenie}\label{cat:prop:nerve-inner-fibration}
Dla dowolnego funktora $F:\mathcal C\to\mathcal D$ małych kategorii morfizm $N(F)$ jest fibracją wewnętrzną.
\end{stwierdzenie}
\begin{proof}
Rozważmy kwadrat podnoszenia dla rogu wewnętrznego. Na mocy twierdzenia~\ref{cat:thm:nerve-inner-horns} róg w $N(\mathcal C)$ ma jednoznaczne wypełnienie $u$. Jego obraz $N(F)(u)$ w $N(\mathcal D)$ ma te same ściany rogu co zadany sympleks $v$ w $N(\mathcal D)$, gdyż kwadrat jest przemienny. Twierdzenie o jednoznaczności wypełnień w $N(\mathcal D)$ daje $N(F)(u)=v$. Zatem $u$ jest wymaganym względnym wypełnieniem.
\end{proof}

\begin{definicja}[Izofibracja kategorii zwykłych]\label{cat:def:ordinary-isofibration}
Funktor $F:\mathcal C\to\mathcal D$ jest izofibracją, jeśli dla każdego obiektu $c\in\mathcal C$ i każdego izomorfizmu $\alpha:F(c)\to d$ w $\mathcal D$ istnieją obiekt $c'\in\mathcal C$ i izomorfizm $a:c\to c'$ z $F(c')=d$ oraz $F(a)=\alpha$.
\end{definicja}
Równoważnie możemy wymagać podnoszenia izomorfizmu $\alpha:d\to F(c)$ z zadanym celem $c$: wystarczy podnieść $\alpha^{-1}$ z zadanym źródłem i wziąć odwrotność otrzymanego izomorfizmu.

\begin{twierdzenie}[Fibracje Kana na nerwach grupoidów]\label{cat:thm:nerve-groupoid-isofibration}
Niech $\mathcal G,\mathcal H$ będą małymi grupoidami i $F:\mathcal G\to\mathcal H$ funktorem. Wtedy $N(F)$ jest fibracją Kana wtedy i tylko wtedy, gdy $F$ jest izofibracją.
\end{twierdzenie}
\begin{proof}
Jeżeli $N(F)$ jest fibracją Kana, to dla $c\in\mathcal G$ i izomorfizmu $\alpha:F(c)\to d$ rozważamy kwadrat z lewym morfizmem $\Lambda^0[1]\to\Delta[1]$. Róg $\Lambda^0[1]$ to wierzchołek $0$, którego obrazem jest $c$, a dolnym sympleksem jest $\alpha$. Podniesienie daje morfizm $a:c\to c'$ z $F(a)=\alpha$ i $F(c')=d$. Morfizm $a$ jest izomorfizmem, gdyż $\mathcal G$ jest grupoidem.

Odwrotnie, załóżmy, że $F$ jest izofibracją. W wymiarze $1$ róg z zadanym źródłem podnosimy bezpośrednio przez definicję izofibracji; róg z zadanym celem podnosimy, stosując ją do odwrotności dolnej krawędzi i odwracając otrzymany morfizm.

Dla $n\geq2$ róg w $N(\mathcal G)$ ma jednoznaczne wypełnienie: dla rogów wewnętrznych wynika to z twierdzenia~\ref{cat:thm:nerve-inner-horns}, a dla zewnętrznych z twierdzenia~\ref{cat:thm:nerve-kan-groupoid} i jego dowodu jednoznaczności w wymiarach co najmniej $2$. Obraz wypełnienia i zadany dolny sympleks w $N(\mathcal H)$ wypełniają ten sam róg, więc są równe przez tę samą jednoznaczność w grupoidzie $\mathcal H$. Otrzymujemy względne wypełnienie dla każdego rogu.
\end{proof}

\begin{przyklad}\label{cat:ex:inner-not-kan-fibration}
Morfizm $N([1])\to\Delta[0]$ jest fibracją wewnętrzną, ale nie jest fibracją Kana, gdyż $N([1])$ nie jest kompleksem Kana (przykład~\ref{cat:ex:nerve-not-kan}). Co więcej, inkluzja grupoidu jednoobiektowego z jedynym morfizmem do grupoidu o dwóch obiektach i dokładnie jednym morfizmie dla każdej uporządkowanej pary obiektów nie jest izofibracją: izomorfizm od obiektu należącego do obrazu do drugiego obiektu nie może zostać podniesiony. Oba nerwy są kompleksami Kana, lecz morfizm między nimi nie jest fibracją Kana.
\end{przyklad}

\subsection{Podnoszenie względem brzegów}

\begin{definicja}[Morfizm z własnością podnoszenia względem brzegów]\label{cat:def:boundary-injective}
Powiemy, że $p:X\to Y$ ma prawą własność podnoszenia względem brzegów, jeśli $\partial\Delta[n]\to\Delta[n]\pitchfork p$ dla wszystkich $n\geq0$. W wymiarze zero brzeg jest pusty, więc warunek wymusza surjektywność $p_0$.
\end{definicja}

\begin{stwierdzenie}\label{cat:prop:boundary-lifting-kan}
Każdy morfizm z prawą własnością podnoszenia względem brzegów jest fibracją Kana.
\end{stwierdzenie}
\begin{proof}
Niech będzie zadany kwadrat podnoszenia dla rogu $\Lambda^k[n]$, $n\geq1$. Oznaczmy górny morfizm przez $a$, dolny przez $b:\Delta[n]\to Y$. Brakująca ściana $\delta^k:\Delta[n-1]\to\Delta[n]$ przecina róg dokładnie w swoim brzegu: dla każdego sympleksu tej ściany jego obraz pomija pewien wierzchołek różny od $k$ wtedy i tylko wtedy, gdy odpowiadający mu morfizm do $[n-1]$ nie jest surjekcją. Dla $n=1$ przecięcie jest puste.

Mamy więc kwadrat podnoszenia z $\partial\Delta[n-1]\to\Delta[n-1]$, którego górnym morfizmem jest ograniczenie $a$ do przecięcia, a dolnym $b\delta^k$. Wybieramy podniesienie $c:\Delta[n-1]\to X$. Zgodność $c$ i $a$ na przecięciu pozwala skleić je w morfizm $a':\partial\Delta[n]\to X$: w każdym stopniu brzeg jest sumą podzbiorów odpowiadających rogowi i brakującej ścianie, zaś naturalność sklejonego morfizmu wynika z naturalności obu ograniczeń. Mamy $pa'=b|_{\partial\Delta[n]}$, ponieważ ta równość zachodzi na obu podzbiorach. Drugi raz stosujemy własność podnoszenia względem brzegu, teraz w wymiarze $n$, uzyskując $l:\Delta[n]\to X$ z $pl=b$ i $l|_{\partial\Delta[n]}=a'$. W szczególności $l|_{\Lambda^k[n]}=a$.
\end{proof}


\subsection{Słabe systemy faktoryzacji}

\begin{definicja}[Słaby system faktoryzacji]\label{cat:def:weak-factorization-system}
Parę klas morfizmów $(\mathcal L,\mathcal R)$ nazywamy słabym systemem faktoryzacji, jeśli
\[
 \mathcal L={}^{\pitchfork}\mathcal R,\qquad
 \mathcal R=\mathcal L^\pitchfork
\]
oraz każdy morfizm $f:A\to B$ ma faktoryzację $f=pi$, gdzie $i:A\to E$ należy do $\mathcal L$, a $p:E\to B$ należy do $\mathcal R$. Nie zakładamy jednoznaczności faktoryzacji ani podniesień.
\end{definicja}

\begin{lemat}[Argument retraktu]\label{cat:lem:factorization-retract}
Niech $f=pi$ będzie faktoryzacją $A\xrightarrow{i}E\xrightarrow{p}B$.
Jeżeli $f\pitchfork p$, to $f$ jest retraktem $i$ w kategorii morfizmów. Jeżeli $i\pitchfork f$, to $f$ jest retraktem $p$.
\end{lemat}
\begin{proof}
W pierwszym przypadku stosujemy własność podnoszenia do kwadratu o lewej krawędzi $f$, prawej $p$, górnej $i$ i dolnej $\id_B$. Otrzymujemy $l:B\to E$ z $lf=i$ i $pl=\id_B$. Pary $(\id_A,l):f\to i$ oraz $(\id_A,p):i\to f$ są morfizmami w kategorii morfizmów, gdyż $i=lf$ oraz $f=pi$. Ich złożenie od $f$ do $f$ ma składowe $\id_A,pl=\id_B$.

W drugim przypadku bierzemy kwadrat o lewej krawędzi $i$, prawej $f$, górnej $\id_A$ i dolnej $p$. Podniesienie $r:E\to A$ spełnia $ri=\id_A$ i $fr=p$. Pary $(i,\id_B):f\to p$ i $(r,\id_B):p\to f$ są morfizmami, a ich złożenie od $f$ do $f$ jest identycznością.
\end{proof}

\begin{stwierdzenie}[Kryterium słabego systemu faktoryzacji]\label{cat:prop:wfs-criterion}
Załóżmy, że $\mathcal L$ i $\mathcal R$ są domknięte względem retraktów, każdy morfizm z $\mathcal L$ ma lewą własność podnoszenia względem każdego morfizmu z $\mathcal R$ i każdy morfizm ma faktoryzację $pi$ z $i\in\mathcal L$, $p\in\mathcal R$. Wówczas $(\mathcal L,\mathcal R)$ jest słabym systemem faktoryzacji.
\end{stwierdzenie}
\begin{proof}
Z założenia $\mathcal L\subseteq{}^\pitchfork\mathcal R$ i $\mathcal R\subseteq\mathcal L^\pitchfork$. Dla $f\in{}^\pitchfork\mathcal R$ wybieramy faktoryzację $f=pi$ jak w założeniu. Mamy $f\pitchfork p$, więc $f$ jest retraktem $i\in\mathcal L$ przez lemat~\ref{cat:lem:factorization-retract}. Stąd $f\in\mathcal L$. Jeżeli natomiast $f\in\mathcal L^\pitchfork$, to $i\pitchfork f$ i ten sam lemat daje, że $f$ jest retraktem $p\in\mathcal R$. Zatem $f\in\mathcal R$.
\end{proof}

\subsection{Argument małego obiektu dla zbiorów symplicjalnych}

Niech $I$ będzie zbiorem monomorfizmów w $\mathbf{sSet}$. W tej części wystarczy przypadek, w którym dziedzina każdego $j:A\to B$ z $I$ ma skończenie wiele sympleksów niezdegenerowanych. Przykładami są rodziny
\[
 I_{\partial}=\{\partial\Delta[n]\hookrightarrow\Delta[n]:n\geq0\},
 \qquad
 I_{\Lambda}=\{\Lambda^k[n]\hookrightarrow\Delta[n]:n\geq1,\ 0\leq k\leq n\}.
\]
W reprezentowalnym zbiorze symplicjalnym $\Delta[n]$ sympleksy niezdegenerowane odpowiadają injekcjom do $[n]$, których jest skończenie wiele. Dla dowolnego podzbioru symplicjalnego $K\subseteq\Delta[n]$ degeneracja sympleksu należącego do $K$ może być rozłożona wewnątrz $K$: jeśli $u=\sigma^*v$ dla surjekcji $\sigma$, wybieramy jej sekcję $\tau$ i mamy $v=\tau^*u\in K$. Zatem niezdegenerowane sympleksy $K$ są dokładnie tymi spośród niezdegenerowanych sympleksów $\Delta[n]$, które należą do $K$.

\begin{lemat}[Małość względem ciągów monomorfizmów]\label{cat:lem:simplicial-finite-smallness}
Jeżeli $A$ ma skończenie wiele sympleksów niezdegenerowanych, a
\[
 Z_0\hookrightarrow Z_1\hookrightarrow Z_2\hookrightarrow\cdots
\]
jest ciągiem monomorfizmów zbiorów symplicjalnych, to każdy morfizm $A\to\operatorname{colim}_{r<\omega}Z_r$ faktoryzuje się przez pewne $Z_r$. Jeżeli dwa morfizmy $A\to Z_r$ mają ten sam obraz w kogranicy, są równe.
\end{lemat}
\begin{proof}
Oznaczmy kogranicę przez $Z$. W każdym stopniu obliczamy ją jako sumę rosnących zbiorów po utożsamieniu $Z_r$ z jego obrazem w $Z$. Włożenia są monomorfizmami, ponieważ przejścia w ciągu są injekcjami w każdym stopniu.

Niech $f:A\to Z$. Dla każdego niezdegenerowanego sympleksu $a$ wybieramy indeks $r_a$ taki, że $f(a)$ należy do obrazu $Z_{r_a}$. Z powodu skończoności zbioru tych sympleksów istnieje jeden indeks $r$ nie mniejszy od wszystkich $r_a$ (dla pustego $A$ bierzemy $r=0$). Każdy sympleks $A$ jest degeneracją sympleksu niezdegenerowanego na mocy lematu~\ref{cat:lem:ez-degeneracy}. Obraz $Z_r$ jest podzbiorem symplicjalnym $Z$, więc zawiera wszystkie $f(a)$, również dla zdegenerowanych $a$. Ograniczając przeciwdziedziny składowych $f$, otrzymujemy morfizm $A\to Z_r$; naturalność wynika z tego, że operatory $Z_r$ są ograniczeniami operatorów $Z$. Ostatnia teza wynika z monomorficzności włożenia $Z_r\to Z$.
\end{proof}

\begin{definicja}[Morfizmy komórkowe względem $I$]\label{cat:def:relative-cell-complex}
Morfizm nazywamy względnym kompleksem komórkowym względem $I$, jeżeli jest złożeniem pozaskończonym zmian kobazy koproduktów morfizmów z $I$. Dokładniej: istnieje ciąg indeksowany liczbami porządkowymi $(Z_\alpha)_{\alpha\leq\lambda}$, z $Z_\beta=\operatorname{colim}_{\alpha<\beta}Z_\alpha$ dla niezerowych granicznych $\beta$, a $Z_\alpha\to Z_{\alpha+1}$ jest zmianą kobazy koproduktu rodziny morfizmów z $I$. Rozważany morfizm to $Z_0\to Z_\lambda$, z dokładnością do izomorfizmu w kategorii morfizmów.
\end{definicja}

\begin{lemat}[Podnoszenie dla kompleksów komórkowych]\label{cat:lem:cell-lifting}
Jeżeli $p\in I^\pitchfork$, to każdy względny kompleks komórkowy względem $I$ ma lewą własność podnoszenia względem $p$.
\end{lemat}
\begin{proof}
Koprodukt morfizmów z $I$ ma lewą własność podnoszenia względem $p$: ograniczamy kwadrat do każdego składnika, wybieramy podniesienia i łączymy je za pomocą własności uniwersalnej koproduktu. Zmiana kobazy zachowuje tę własność na mocy stwierdzenia~\ref{cat:prop:lifting-base-change}.

Dla kwadratu z lewym morfizmem $Z_0\to Z_\lambda$, górnym $a:Z_0\to E$ i dolnym $b:Z_\lambda\to Y$, budujemy zgodne morfizmy $l_\alpha:Z_\alpha\to E$ z $pl_\alpha=b|_{Z_\alpha}$. Zaczynamy od $l_0=a$. Na etapie następnikowym stosujemy własność podnoszenia $Z_\alpha\to Z_{\alpha+1}$ do kwadratu z $l_\alpha$ i $b|_{Z_{\alpha+1}}$. Na etapie granicznym korzystamy z własności uniwersalnej kogranicy zgodnych $l_\alpha$; równość z $b$ po złożeniu z $p$ zachodzi na wszystkich składnikach, zatem zachodzi na kogranicy. Rekursja pozaskończona z wyborem podniesień daje $l_\lambda$, czyli żądany morfizm. W przypadku pustego ciągu lewy morfizm jest izomorfizmem.
\end{proof}

\begin{twierdzenie}[Argument małego obiektu w wersji przeliczalnej]\label{cat:thm:small-object-sset}
Niech $I$ będzie zbiorem monomorfizmów zbiorów symplicjalnych, których dziedziny mają skończenie wiele sympleksów niezdegenerowanych. Każdy morfizm $f:X\to Y$ ma faktoryzację $f=pi$, w której $i$ jest względnym kompleksem komórkowym względem $I$ (wystarcza ciąg długości $\omega$), a $p\in I^\pitchfork$. Para
\[
 ({}^\pitchfork(I^\pitchfork),I^\pitchfork)
\]
jest słabym systemem faktoryzacji. Każdy morfizm z lewej klasy jest retraktem względnego kompleksu komórkowego względem $I$.
\end{twierdzenie}
\begin{proof}
Zaczynamy od $X_0=X$ oraz $f_0=f$. Dla każdego $r\geq0$ tworzymy zbiór $S_r$ wszystkich trójek $(j,a,b)$, gdzie $j:A_j\to B_j$ należy do $I$, $a:A_j\to X_r$, $b:B_j\to Y$ oraz $f_ra=bj$. Jest to zbiór, ponieważ $I$ jest zbiorem, a kategoria zbiorów symplicjalnych jest lokalnie mała. Tworzymy sumę amalgamowaną
\[
 X_{r+1}=X_r\amalg_{\coprod_{(j,a,b)\in S_r}A_j}
                       \coprod_{(j,a,b)\in S_r}B_j,
\]
gdzie górne morfizmy do $X_r$ są dane przez $a$, a drugi morfizm z koproduktu dziedzin przez $j$ na każdym składniku. Własność uniwersalna tej sumy amalgamowanej daje $f_{r+1}:X_{r+1}\to Y$, którego ograniczenie do $X_r$ to $f_r$, a do każdego $B_j$ to $b$. Zatem każdy kwadrat z $S_r$ ma podniesienie po przejściu do $X_{r+1}$.

\begin{center}
\begin{tikzpicture}[>=stealth]
 \node (A) at (0,2.4) {$\coprod_{S_r}A_j$};
 \node (X) at (5,2.4) {$X_r$};
 \node (B) at (0,0) {$\coprod_{S_r}B_j$};
 \node (Z) at (5,0) {$X_{r+1}$};
 \node (Y) at (8,-1.7) {$Y$};
 \draw[->] (A)--node[above] {$\coprod a$}(X);
 \draw[->,hook] (A)--node[left] {$\coprod j$}(B);
 \draw[->,hook] (X)--(Z);
 \draw[->] (B)--(Z);
 \draw[->] (Z)--node[above,sloped] {$f_{r+1}$}(Y);
 \draw[->] (X) to[bend left=22] node[right] {$f_r$}(Y);
 \draw[->] (B) to[bend right=15] node[below] {$\coprod b$}(Y);
 \node at (2.5,1.1) {$\mathrm{PO}$};
\end{tikzpicture}
\end{center}

Każdy $X_r\to X_{r+1}$ jest monomorfizmem: koprodukt monomorfizmów jest stopniowo injekcją, a zmiana kobazy injekcji zbiorów $U\hookrightarrow V$ przez funkcję $U\to W$ jest injekcją $W\hookrightarrow W\amalg_U V$. Istotnie, po utożsamieniu $U$ z podzbiorem $V$ zbiór po prawej jest izomorficzny z rozłączną sumą $W\amalg(V\setminus U)$, gdzie elementy $U$ przechodzą na swoje obrazy w $W$. Stosujemy ten rachunek w każdym stopniu.

Niech $X_\omega=\operatorname{colim}_{r<\omega}X_r$, a $p:X_\omega\to Y$ będzie indukowany przez zgodne $f_r$. Oznaczamy przez $i:X\to X_\omega$ morfizm kanoniczny. Mamy $pi=f$, a $i$ jest względnym kompleksem komórkowym. Sprawdzimy $p\in I^\pitchfork$. Dla dowolnego kwadratu $pa=bj$ z $j:A_j\to B_j$ w $I$ i $a:A_j\to X_\omega$ lemat~\ref{cat:lem:simplicial-finite-smallness} daje faktoryzację $a$ przez $a_r:A_j\to X_r$. Ponieważ $p|_{X_r}=f_r$, mamy $f_ra_r=bj$, a więc $(j,a_r,b)\in S_r$. Morfizm $B_j\to X_{r+1}\to X_\omega$ jest podniesieniem wyjściowego kwadratu.

Z lematu~\ref{cat:lem:cell-lifting} mamy $i\in{}^\pitchfork(I^\pitchfork)$. Obie klasy są domknięte względem retraktów na mocy stwierdzenia~\ref{cat:prop:lifting-retract} i mają wzajemną własność podnoszenia. Zatem stwierdzenie~\ref{cat:prop:wfs-criterion} daje słaby system faktoryzacji. Jeżeli $f$ należy do lewej klasy, to dla skonstruowanej faktoryzacji $f\pitchfork p$, więc lemat~\ref{cat:lem:factorization-retract} daje, że $f$ jest retraktem $i$.
\end{proof}

\begin{wniosek}[Systemy generowane przez rogi]\label{cat:cor:horn-wfs}
Fibracje Kana, fibracje wewnętrzne, lewe i prawe fibracje są prawymi klasami słabych systemów faktoryzacji generowanych odpowiednio przez wszystkie rogi, rogi wewnętrzne, rogi z $k<n$ oraz rogi z $k>0$.
\end{wniosek}
\begin{proof}
Każda z wymienionych rodzin jest zbiorem monomorfizmów o dziedzinach mających skończenie wiele niezdegenerowanych sympleksów. Stosujemy twierdzenie~\ref{cat:thm:small-object-sset} i definicję~\ref{cat:def:simplicial-fibrations}.
\end{proof}

\begin{definicja}[Morfizmy anodyniczne]\label{cat:def:anodyne}
Morfizm zbiorów symplicjalnych nazywamy anodynicznym, jeśli ma lewą własność podnoszenia względem wszystkich fibracji Kana. Wewnętrznie anodyniczny morfizm ma lewą własność podnoszenia względem wszystkich fibracji wewnętrznych. Lewostronnie i prawostronnie anodyniczne morfizmy definiujemy analogicznie względem lewych i prawych fibracji.
\end{definicja}
Każdy z tych morfizmów jest retraktem odpowiedniego względnego kompleksu komórkowego, a więc monomorfizmem: morfizmy przejścia są stopniowo injekcjami, złożenie pozaskończone injekcji jest injekcją i retrakt monomorfizmu jest monomorfizmem. Dla ostatniego faktu, jeśli $f$ jest retraktem $g$ z kwadratami jak w definicji~\ref{cat:def:arrow-retract} i $fa=fb$, to $gua=vfa=vfb=gub$, skąd $ua=ub$ i $a=b$ po złożeniu z $r$.

\subsection{Aksjomaty kategorii modelowych}

\begin{definicja}[Kategoria modelowa]\label{cat:def:model-category}
Kategorią modelową nazywamy kategorię $\mathcal M$ mającą wszystkie małe granice i kogranice, wyposażoną w klasy morfizmów $\mathcal W$, $\mathcal C$, $\mathcal F$, nazywane odpowiednio słabymi równoważnościami, kofibracjami i fibracjami, spełniające następujące warunki:
\begin{enumerate}
 \item każda z trzech klas zawiera izomorfizmy i jest domknięta względem retraktów;
 \item jeżeli dla składanych morfizmów $f,g$ dwa spośród $f,g,gf$ należą do $\mathcal W$, to trzeci również należy do $\mathcal W$ (własność dwóch z trzech);
 \item każda kofibracja ma lewą własność podnoszenia względem fibracji należących do $\mathcal W$, a każda kofibracja należąca do $\mathcal W$ ma lewą własność podnoszenia względem każdej fibracji;
 \item każdy morfizm $f$ ma obie faktoryzacje
 \[
 f=pi=qj,\quad
 i\in\mathcal C,\ p\in\mathcal F\cap\mathcal W,\quad
 j\in\mathcal C\cap\mathcal W,\ q\in\mathcal F.
 \]
\end{enumerate}
W tej konwencji nie wymagamy wyboru funktorialnych faktoryzacji. Kofibracje i fibracje należące do $\mathcal W$ nazywamy odpowiednio trywialnymi kofibracjami i trywialnymi fibracjami.
\end{definicja}

\begin{stwierdzenie}[Dwa słabe systemy faktoryzacji]\label{cat:prop:model-wfs}
W kategorii modelowej parami słabych systemów faktoryzacji są
\[
 (\mathcal C,\mathcal F\cap\mathcal W),\qquad
 (\mathcal C\cap\mathcal W,\mathcal F).
\]
Odwrotnie, jeśli w kategorii z małymi granicami i kogranicami $\mathcal W$ zawiera izomorfizmy, jest domknięta względem retraktów i spełnia własność dwóch z trzech, a obie wskazane pary są słabymi systemami faktoryzacji, to otrzymujemy kategorię modelową.
\end{stwierdzenie}
\begin{proof}
Dla pierwszego kierunku w obu parach klasy są domknięte względem retraktów, aksjomat podnoszenia daje wzajemną własność podnoszenia, a aksjomat faktoryzacji daje rozkład każdego morfizmu. Stosujemy stwierdzenie~\ref{cat:prop:wfs-criterion}. Dla drugiego kierunku klasy $\mathcal C$ i $\mathcal F$ zawierają izomorfizmy i są domknięte względem retraktów na mocy stwierdzeń~\ref{cat:prop:lifting-composition} i~\ref{cat:prop:lifting-retract}, ponieważ są klasami określonymi przez własność podnoszenia. Pozostałe warunki wynikają wprost z założeń o $\mathcal W$ i obu systemach.
\end{proof}

\begin{definicja}[Obiekty kofibrantne i fibrantne]\label{cat:def:cofibrant-fibrant}
Niech $0$ i $1$ oznaczają odpowiednio obiekt początkowy i końcowy kategorii modelowej. Obiekt $X$ jest kofibrantny, jeśli $0\to X$ jest kofibracją, i fibrantny, jeśli $X\to1$ jest fibracją. Obiekt zarazem kofibrantny i fibrantny nazywamy bifibrantnym.
\end{definicja}

\begin{wniosek}[Zastąpienia]\label{cat:cor:model-replacements}
Każdy obiekt $X$ ma trywialną fibrację $QX\to X$ z obiektu kofibrantnego oraz trywialną kofibrację $X\to RX$ do obiektu fibrantnego. Jeżeli $X$ jest kofibrantny, to tak skonstruowany $RX$ jest bifibrantny. Jeżeli $X$ jest fibrantny, to $QX$ jest bifibrantny.
\end{wniosek}
\begin{proof}
Faktoryzujemy $0\to X$ przez kofibrację $0\to QX$ i trywialną fibrację $QX\to X$, a $X\to1$ przez trywialną kofibrację $X\to RX$ i fibrację $RX\to1$. Jeżeli $0\to X$ jest kofibracją, to złożenie $0\to X\to RX$ jest kofibracją, ponieważ lewa klasa pierwszego systemu jest domknięta względem składania. Jeżeli $X\to1$ jest fibracją, to złożenie $QX\to X\to1$ jest fibracją przez domknięcie prawej klasy drugiego systemu. W szczególności każdy obiekt jest połączony z obiektem bifibrantnym zygzakiem słabych równoważności $X\leftarrow QX\to RQX$.
\end{proof}

\begin{przyklad}[Dyskretna struktura modelowa]\label{cat:ex:discrete-model-structure}
W każdej kategorii z małymi granicami i kogranicami można przyjąć za $\mathcal W$ izomorfizmy, a za $\mathcal C$ i $\mathcal F$ wszystkie morfizmy. Izomorfizmy spełniają własność dwóch z trzech: jeżeli $f$ i $gf$ są izomorfizmami, to $g=(gf)f^{-1}$ jest izomorfizmem, a jeżeli $g$ i $gf$ są izomorfizmami, to $f=g^{-1}(gf)$ jest izomorfizmem. Retrakt izomorfizmu $g:C\to D$ przez dane $u,v,r,s$ z definicji~\ref{cat:def:arrow-retract} jest izomorfizmem: odwrotnością $f$ jest $rg^{-1}v$, gdyż $rg^{-1}vf=rg^{-1}gu=ru=\id$ oraz $frg^{-1}v=sgg^{-1}v=sv=\id$. Własności podnoszenia wynikają ze stwierdzenia~\ref{cat:prop:lifting-composition}, a faktoryzacje to $f=\id_B f=f\id_A$ dla $f:A\to B$. Wszystkie obiekty są bifibrantne.
\end{przyklad}
